\section{Lifting actual tangent-symbol regularity from the closed fibre}

Let $B$ be a commutative local ring with residue field $k$, let $R$ be
a commutative $B$-algebra, let $I,N\subset R$ be entire ideals and
$v\in N$. Write $D_q=\operatorname{im}(N^q\to R/(I+N^{q+1}))$, with
the scalar action induced from the given $B$-algebra structure. The
actual multiplication map $f_q:D_q\to D_{q+1}$ sends the original
numerator $b$ to $vb$. Assume each $D_q$ is finite free over $B$ and
each actual residue-field map $k\otimes_B f_q$ is injective. The latter
is an explicit fibre condition; no conclusion about a particular
geometric family is asserted without proving that condition there.

\begin{theorem}[Closed-fibre injectivity splits the actual symbol map]
\label{hir:local-actual-symbol-map-splits}
\uses{hir:strict-parameter-symbol-map}
\lean{Hironaka.HironakaLocalSymbolRegularity.adicSymbolMul_split_of_residue_injective}
\leanok
For every $q$, the actual $B$-linear map $f_q$ has a $B$-linear left
inverse, and is therefore injective. It suffices here that $D_q$ be
finite and $D_{q+1}$ finite free.
\end{theorem}
\begin{proof}
The local-ring split-injection theorem applies to the genuine map
$f_q$, since its source is finite, its target is finite free and its
residue-field tensor is injective. It supplies a $B$-linear map $r_q$
with $r_qf_q=1$. If $f_qx=f_qy$, applying $r_q$ proves $x=y$.
This concerns the actual layer elements and their unchanged scalar
structure, rather than only a comparison of dimensions.
\end{proof}

\begin{theorem}[The original whole-ideal regular-symbol implication follows]
\label{hir:local-actual-symbol-whole-regularity}
\uses{hir:local-actual-symbol-map-splits}
\lean{Hironaka.HironakaLocalSymbolRegularity.regular_symbol_of_residue_injective}
\leanok
For every original $b\in N^q$,
\[
 vb\in I+N^{q+2}\quad\Longrightarrow\quad b\in I+N^{q+1}.
\]
Thus the ordinary entire-ideal regular-symbol hypothesis is internally
proved from the preceding finite-module and actual fibre assumptions.
\end{theorem}
\begin{proof}
The original $b$ represents an element $[b]\in D_q$. The given
membership says precisely that the genuine class $[vb]=f_q[b]$ in
$D_{q+1}$ is zero. Since $f_q$ is injective, $[b]=0$ in $D_q$.
The original ideal quotient dictionary now gives
$b\in I+N^{q+1}$. Every numerator and denominator is unchanged.
\end{proof}

\begin{theorem}[Each literal parameter-quotient layer remains finite free]
\label{hir:local-actual-symbol-quotient-layers-free}
\uses{hir:local-actual-symbol-whole-regularity,hir:strict-parameter-whole-layer-kernel,hir:strict-parameter-degree-zero}
\lean{Hironaka.HironakaLocalSymbolRegularity.adicLayer_after_parameter_finite_free}
\leanok
For every $q$, the literal layer
$\operatorname{im}(N^q\to R/(I+(v)+N^{q+1}))$ is finite free over the
original $B$.
\end{theorem}
\begin{proof}
The canonical original layer map onto this layer is surjective, so
finiteness descends from $D_q$. At degree zero it is also injective,
since $v\in N$ leaves the degree-zero denominator unchanged; hence
it transfers the given free structure on $D_0$.

In positive degree $q+1$, the whole-ideal regularity proved above and
the strict filtered quotient theorem identify the entire kernel of
the canonical layer map with the image of the actual $f_q$. This
equality remains true after restricting scalars from $R$ to $B$:
both kernels and images have the same underlying elements and the
original $B$-action. We therefore have the genuine exact sequence
$D_q\xrightarrow{f_q}D_{q+1}\to D_{q+1}(I+(v))\to0$.
The local-ring presentation theorem says its cokernel is free when
the first two modules are finite, the middle one free and the
residue-field tensor of the first map injective. These are precisely
the stated assumptions. Its conclusion concerns the literal target
layer in this sequence, completing every positive degree. No
separatedness, generic localization or additional basis is used.
\end{proof}
